%% kapa: exam=1 round=23 year=2012 period=2 form=A subject=appraisal_law subject_ko=부동산관계법규 seq=22 number=62
%% source: Q-Net 감정평가사 기출문제 (cst00309) artlSeq=5052580 file=raw/1008/5052580/2012년 제23회 감정평가사 제1차시험 기출문제(2교시 A형).pdf
%% conversion: kordoc 4.24.51 -> tools/gen1.py (auto); review=proofread-sonnet; flags=none
\begin{question}{62}
부동산 가격공시 및 감정평가에 관한 법령상 개별공시지가의 결정ㆍ공시에 관한 설명으로 옳지 않은 것은? (공시기준일은 1월 1일로 함)
\begin{choices}
  \item 개별공시지가를 결정ㆍ공시하지 아니한 표준지에 대하여는 당해 토지의 공시지가를 개별공시지가로 본다.
  \item 시장ㆍ군수 또는 구청장은 조세의 부과대상이 아닌 토지에 대해서는 개별공시지가를 결정ㆍ공시하지 아니할 수 있다.
  \item 2011년 11월 15일 분할ㆍ합병된 토지에 대해서는 2012년 5월 31일까지 개별공시지가를 결정ㆍ공시하여야 한다.
  \item 2012년 5월 15일 토지의 형질변경으로 지목이 변경된 토지에 대해서는 2012년 7월 1일을 기준일로 하여 개별공시지가를 결정ㆍ공시하여야 한다.
  \item 시장ㆍ군수 또는 구청장이 개별공시지가를 결정ㆍ공시하기 위하여 개별토지의 가격을 산정한 때에는 토지소유자 등의 의견을 들어야 하나, 미리 관계 중앙행정기관의 장과 협의한 경우에는 의견청취를 생략할 수 있다.
\end{choices}
\end{question}
\answer{5}
